\documentclass[10pt,aspectratio=169]{beamer}

\usetheme{metropolis}
\usepackage{appendixnumberbeamer}
\usepackage{kotex}
\usepackage{pgfpages}

% Note mode switch:
% - default: hide notes (for make slide)
% - define \shownotes to enable note pages (for make_note)
\newif\ifshownotes
\ifdefined\shownotes
  \shownotestrue
\else
  \shownotesfalse
\fi

\ifshownotes
  \setbeameroption{show notes}
  \setbeamertemplate{note page}[plain]
  \pgfpagesuselayout{2 on 1}[a4paper,border shrink=5mm]
\else
  \setbeameroption{hide notes}
\fi

\title{Metropolis Beamer Demo}
\subtitle{Lecture 01}
\author{Your Name}
\institute{Your Organization}
\date{\today}

\begin{document}

\maketitle

\begin{frame}{Agenda}
  \tableofcontents
  \note{발표의 전체 흐름을 간단히 소개합니다.}
\end{frame}

\section{Introduction}
\begin{frame}{Why Beamer + Metropolis?}
  \begin{itemize}
    \item Clean and modern style
    \item Fast PDF workflow
    \item Easy version control
  \end{itemize}
  \note{Metropolis의 장점 3가지를 강조해 주세ㅇㅇㅇ요.}
\end{frame}

\section{Content}
\begin{frame}{Key Message}
  \begin{block}{Takeaway}
    Keep slides simple and focused.
  \end{block}
  \note{핵심 문장을 먼저 말하고 보충 설명을 dddxxㅊㅊㅊxx이어갑니다.}
\end{frame}

\begin{frame}{Two-column Example}
  \begin{columns}[T,onlytextwidth]
    \column{0.5\textwidth}
      \textbf{Left}
      \begin{itemize}
        \item Point A
        \item Point B
      \end{itemize}
    \column{0.5\textwidth}
      \textbf{Right}
      \begin{enumerate}
        \item Step 1ㄴㄴㄴㄴㄴ
        \item Step 2
      \end{enumerate}
  \end{columns}
  \note{좌/우 비교 구조를 설명할 때 시선 이동을 ㅊㅊㅊ유도합니다.}
\end{frame}

\section{Summary}
\begin{frame}{Summary}
  \begin{itemize}
    \item Use Metropolis theme
    \item Build with latexmk via Makefile
    \item Speaker notes via \texttt{make\_note}
  \end{itemize}
  \note{마지막에 빌드 명령을 다시 한 번 안내합니다.}
\end{frame}

\begin{frame}[standout]
  Thank you!
  \note{질의응답으로 자연스럽게 전환합니다.}
\end{frame}

\end{document}
